По одному из билетов, $K$ -- линейное пространство над $\mathbb{F}_p$, $\mathbb{F}_{p^n}$ -- линейное пространство над $K$. Тогда $|K| = p^m, p^n = (p^m)^k$ для некоторых натуральных $m, k$. 

Почему такое поле существует для любого $m \leq n$: это следует из того, что $(x^{p^m} - x)|(x^{p^n} - x)$ для $m|n$, поэтому можно взять корни делителя, и они будут образовывать подполе нужного размера.